\documentclass[11pt]{article}
\pdftrailerid{}
\pdfinfoomitdate=1
\pdfinfo{/CreationDate (D:20260924000000Z) /ModDate (D:20260924000000Z)}
\usepackage[T1]{fontenc}
\usepackage{lmodern}
\usepackage[margin=1.08in]{geometry}
\usepackage{amsmath,amssymb,amsthm,mathtools}
\usepackage{microtype}
\usepackage{aliascnt}
\usepackage{enumitem}
\usepackage{booktabs}
\usepackage{xcolor}
\usepackage{xurl}
\usepackage{tikz}
\usetikzlibrary{arrows.meta,positioning}
\usepackage[colorlinks=true,linkcolor=blue!45!black,citecolor=blue!45!black,urlcolor=blue!45!black]{hyperref}
\usepackage[capitalise,nameinlink,noabbrev]{cleveref}
\pdfgentounicode=1
\pdfglyphtounicode{tfm:lmsy10/mapsto}{007C}
\pdfglyphtounicode{tfm:lmex10/parenleftbig}{0028}
\pdfglyphtounicode{tfm:lmex10/parenrightbig}{0029}
\pdfglyphtounicode{tfm:lmex10/bracketleftbig}{005B}
\pdfglyphtounicode{tfm:lmex10/bracketrightbig}{005D}
\pdfglyphtounicode{tfm:lmex10/braceleftbig}{007B}
\pdfglyphtounicode{tfm:lmex10/bracerightbig}{007D}
\pdfglyphtounicode{tfm:lmex10/vextendsingle}{23D0}
\pdfglyphtounicode{tfm:lmex10/parenleftbigg}{0028}
\pdfglyphtounicode{tfm:lmex10/parenrightbigg}{0029}
\pdfglyphtounicode{tfm:lmex10/bracketleftbigg}{005B}
\pdfglyphtounicode{tfm:lmex10/bracketrightbigg}{005D}
\pdfglyphtounicode{tfm:lmex10/floorleftbigg}{230A}
\pdfglyphtounicode{tfm:lmex10/floorrightbigg}{230B}
\pdfglyphtounicode{tfm:lmex10/braceleftbigg}{007B}
\pdfglyphtounicode{tfm:lmex10/bracerightbigg}{007D}
\pdfglyphtounicode{tfm:lmex10/angbracketleftbigg}{27E8}
\pdfglyphtounicode{tfm:lmex10/angbracketrightbigg}{27E9}
\pdfglyphtounicode{tfm:lmex10/parenleftBigg}{0028}
\pdfglyphtounicode{tfm:lmex10/parenrightBigg}{0029}
\pdfglyphtounicode{tfm:lmex10/parenlefttp}{239B}
\pdfglyphtounicode{tfm:lmex10/parenrighttp}{239E}
\pdfglyphtounicode{tfm:lmex10/bracelefttp}{23A7}
\pdfglyphtounicode{tfm:lmex10/bracerighttp}{23AB}
\pdfglyphtounicode{tfm:lmex10/braceleftbt}{23A9}
\pdfglyphtounicode{tfm:lmex10/bracerightbt}{23AD}
\pdfglyphtounicode{tfm:lmex10/braceleftmid}{23A8}
\pdfglyphtounicode{tfm:lmex10/bracerightmid}{23AC}
\pdfglyphtounicode{tfm:lmex10/parenleftbt}{239D}
\pdfglyphtounicode{tfm:lmex10/parenrightbt}{23A0}
\pdfglyphtounicode{tfm:lmex10/circleplusdisplay}{2A01}
\pdfglyphtounicode{tfm:lmex10/summationtext}{2211}
\pdfglyphtounicode{tfm:lmex10/integraltext}{222B}
\pdfglyphtounicode{tfm:lmex10/summationdisplay}{2211}
\pdfglyphtounicode{tfm:lmex10/productdisplay}{220F}
\pdfglyphtounicode{tfm:lmex10/integraldisplay}{222B}
\pdfglyphtounicode{tfm:lmex10/tildewide}{0303}
\pdfglyphtounicode{tfm:lmex10/radicalbig}{221A}
\pdfglyphtounicode{tfm:lmex10/radicalBig}{221A}
\pdfglyphtounicode{tfm:lmex10/radicalbigg}{221A}
\pdfglyphtounicode{tfm:lmex10/radicalBigg}{221A}
\numberwithin{equation}{section}
\newtheorem{theorem}{Theorem}[section]
\newaliascnt{lemma}{theorem}
\newtheorem{lemma}[lemma]{Lemma}
\aliascntresetthe{lemma}
\Crefname{lemma}{Lemma}{Lemmas}
\newaliascnt{proposition}{theorem}
\newtheorem{proposition}[proposition]{Proposition}
\aliascntresetthe{proposition}
\Crefname{proposition}{Proposition}{Propositions}
\newaliascnt{corollary}{theorem}
\newtheorem{corollary}[corollary]{Corollary}
\aliascntresetthe{corollary}
\Crefname{corollary}{Corollary}{Corollaries}
\theoremstyle{definition}
\newaliascnt{definition}{theorem}
\newtheorem{definition}[definition]{Definition}
\aliascntresetthe{definition}
\Crefname{definition}{Definition}{Definitions}
\theoremstyle{remark}
\newaliascnt{remark}{theorem}
\newtheorem{remark}[remark]{Remark}
\aliascntresetthe{remark}
\Crefname{remark}{Remark}{Remarks}
\newcommand{\R}{\mathbb R}
\newcommand{\Hh}{\mathcal H}
\newcommand{\E}{\mathbb E}
\newcommand{\op}{\mathrm{op}}
\DeclareMathOperator{\GL}{GL}
\DeclareMathOperator{\OPT}{OPT}
\DeclareMathOperator{\TV}{TV}
\setlist[enumerate]{itemsep=3pt,topsep=5pt}
\setlist[itemize]{itemsep=3pt,topsep=5pt}
\allowdisplaybreaks[1]
\hypersetup{
  pdftitle={Near-square-root logarithmic integrality gaps for uniform sparsest cut},
  pdfauthor={OpenAI},
  pdfsubject={Uniform-demand Goemans--Linial semidefinite integrality gaps},
  pdfkeywords={uniform sparsest cut, integrality gap, negative type, L1 embeddings}}

\title{Near-square-root logarithmic integrality gaps\\for uniform sparsest cut}
\author{OpenAI}
\date{September 24, 2026}

\begin{document}
\maketitle
\begin{abstract}
We construct uniform sparsest-cut instances whose Goemans--Linial
semidefinite integrality gap is at least
$c\sqrt{\log n}/(\log\log n)^3$ along a sequence $n\to\infty$.
The demand is one between every pair of distinct vertices, and the
capacities are nonnegative real numbers. This matches the
Arora--Rao--Vazirani upper bound up to a power of $\log\log n$.
\end{abstract}

\section{Introduction}
\label{sec:introduction}

Sparsest cut asks for a partition with small crossing capacity relative
to the demand it separates. It is a basic graph-partitioning problem,
and its relaxations connect approximation algorithms with the geometry
of finite metric spaces. When every pair of vertices has the same
demand, the denominator depends only on the two part sizes. An
integrality-gap lower bound in this uniform case requires an obstruction
to preserving the average distance over all pairs.

Let $C=(c_{ij})$ be symmetric nonnegative capacities on the unordered
pairs of distinct vertices in $[n]=\{1,\ldots,n\}$. The uniform
sparsest-cut optimum is
\begin{equation}\label{eq:opt}
\OPT(C)=\min_{\varnothing\ne B\subsetneq[n]}
\frac{\sum_{i\in B,\,j\notin B}c_{ij}}{|B|(n-|B|)}.
\end{equation}
Thus every pair has demand exactly one. We allow arbitrary capacities;
no regularity or degree-weighted vertex measure is assumed.

A \emph{negative-type semimetric} on a finite set is a function
$d(i,j)=|x_i-x_j|^2$, for vectors in a real Hilbert space, that satisfies
all triangle inequalities. Distinct points may have distance zero.
The Goemans--Linial relaxation~\cite{Goe97,Lin02} is
\begin{equation}\label{eq:gl}
\GL(C)=\min\left\{
\sum_{i<j}c_{ij}d(i,j):\
\sum_{i<j}d(i,j)=1,\quad d\text{ is a negative-type semimetric}
\right\}.
\end{equation}
Every finite Hilbert configuration can be realized in $\R^n$, so this
is the usual squared-Euclidean SDP with triangle inequalities.
There is no unit-norm constraint on the vectors. Cut distances are
feasible after normalization, and hence $\GL(C)\le\OPT(C)$.

\begin{theorem}\label{thm:main}
There are an absolute constant $c>0$ and capacity instances $C^{(j)}$
on $n_j\to\infty$ vertices, with $\GL(C^{(j)})>0$, such that
\[
\frac{\OPT(C^{(j)})}{\GL(C^{(j)})}
\ge c\,\frac{\sqrt{\log n_j}}{(\log\log n_j)^3}
\]
for all sufficiently large $j$. The demand is one between every pair
of distinct vertices.
\end{theorem}

\subsection{History and the uniform-demand obstacle}

The metric viewpoint on cuts was developed by Linial, London, and
Rabinovich~\cite{LLR95}: finite $\ell_1$ distances are nonnegative
combinations of cut distances, and an embedding into $\ell_1$ can be
rounded to a cut. For uniform demands, the relevant quantity is the
average distance preserved by a nonexpanding map. Rabinovich developed
this average-distortion framework and its connection to uniform
sparsest cut~\cite{Rab03,Rab08}. Here a map is \emph{nonexpanding}, or a
\emph{contraction}, if its distance on each pair is at most the original
distance. The proof below constructs a negative-type semimetric for
which every contraction into $\ell_1$ preserves only a small fraction
of the average distance.

Arora, Rao, and Vazirani gave an $O(\sqrt{\log n})$ approximation for
uniform sparsest cut through their semidefinite relaxation and geometric
separation theorem~\cite{ARV09}. Their work, first presented in 2004,
provides the upper-bound benchmark for \Cref{thm:main}. Devanur, Khot,
Saket, and Vishnoi obtained an $\Omega(\log\log n)$ integrality gap,
refuting the uniform constant-gap conjecture~\cite{DKSV06}.
Kane and Meka subsequently proved the stronger bound
$\exp(\Omega(\sqrt{\log\log n}))$ through pseudorandom generators for
Lipschitz functions of polynomials~\cite{KM13}. Their work also gives
lower bounds for strengthened relaxations; the present paper concerns
the basic Goemans--Linial SDP. Our bound reaches the exponent $1/2$ in
$\log n$, with the displayed iterated-logarithmic loss, and quantitatively
strengthens the earlier negative resolution of the uniform constant-gap
conjecture.

For arbitrary demands, the Goemans--Linial Conjecture predicted a constant
integrality gap, equivalently a uniform bound on the $\ell_1$ distortion
of finite negative-type metrics. Khot and Vishnoi disproved this
conjecture~\cite{KV15}. A geometric line of work based on the Heisenberg
group was developed by Lee and Naor~\cite{LN06}, Cheeger and
Kleiner~\cite{CK10}, and Cheeger, Kleiner, and Naor~\cite{CKN09,CKN11}. Naor and Young obtained the sharp general-demand
lower bound $\Omega(\sqrt{\log n})$~\cite{NY18}; Chang, Naor, and Ren
proved the matching upper bound~\cite{CNR25STOC,CNR25}.
These lower bounds do not imply the uniform-demand conclusion.
Indeed, Cheeger, Kleiner, and Naor explain explicitly that their
bounded-doubling examples have constant average distortion into a line,
so their construction cannot yield a diverging uniform gap
\cite[Remark~1.1]{CKN09}. This distinction between worst-pair distortion
and average-distance preservation is the obstacle addressed here.

\subsection{Proof strategy}

The construction has two tasks: build a negative-type semimetric with
large average distance, and force every $\ell_1$ contraction to have
small average distance for the same probability measure. The final
passage to capacities uses cut-cone duality in the average-distortion
framework of Linial--London--Rabinovich and Rabinovich
\cite{LLR95,Rab03,Rab08}.
We give the full finite-dimensional argument in \Cref{sec:reduction},
including the passage to exactly uniform demands and positivity of the
SDP optimum.

\paragraph{A common parameter space, many rounded charts.}
Fix a large integer $m$ and the parameter cube $U=(-2,2)^m$.
For each $s\in\{1,\ldots,S\}$, choose a linear map
$J_s:\R^m\to\R^N$ with rows $u_{s,i}^{T}$, and round its coordinates
to a lattice of mesh $\tau>0$. The resulting vertex label is
\[
 v_s(\theta)=(s,X_s(\theta)),\qquad
 X_s(\theta)_i=\tau\left\lfloor
 \frac{\langle u_{s,i},\theta\rangle}{\tau}\right\rfloor.
\]
Each chart partitions the same cube into cells cut out by hyperplanes;
$V$ consists of all labels attained off those hyperplanes.
The directions are fixed before choosing any contraction. Their
simultaneous guarantee, proved in \Cref{dir:section} and
Appendix~\ref{app:directions}, ensures that every unit direction has
small projections onto all normals in at least half the charts.
The powers of $m$ determining $S,N,\tau$ are chosen so that
$|V|\le\exp(Cm\log m)$.

\paragraph{Large global distances and small interface costs.}
\Cref{ker:sec} builds Hilbert vectors for the labels using one common
smooth feature map of $\theta$. Squared distances between these features
approximate $c_*\sqrt m\,|\theta-\eta|$ with a uniform additive error,
where $c_*>0$ is absolute.
Small additional Fourier components allow extensions away from each
chart subspace whose coordinate derivatives have small inner products
against all feature vectors. These derivative profiles control the
cost of changing one rounded coordinate. The extension operators may
have large norms; sufficiently fine rounding controls their effect on
the feature values.

The resulting squared distances need not yet satisfy every triangle
inequality. \Cref{sec:metric} repairs them using a scale integral
of Gaussian positive-definite kernels, building on the classical
Gaussian-kernel principle of Schoenberg~\cite{Sch38}. The local angle
estimate is indispensable: a bounded global error alone would overwhelm
the tiny costs at cell interfaces. Together, the local and global angle
bounds permit a repair of bounded total size with only a logarithmic
loss at those interfaces. In the resulting semimetric $d$, with
tolerance $r=m^{-8}$, labels from different charts at
the same parameter have distance at most $Cr$, while adjacent labels
within a chart have the much smaller bound in
\eqref{met:eq:adjacent}. Independent uniform parameters in
$T=(-1,1)^m$ have expected first-chart distance at least $cm$.

\paragraph{Simultaneous constraints on a contraction.}
Given any contraction $F:V\to\ell_1^D$, define
$f_s(\theta)=F(v_s(\theta))$ almost everywhere on $U$.
The same-parameter estimate makes these functions uniformly close.
After smoothing, their gradients are close as well. In each chart,
small jumps across cell interfaces express each scalar gradient as a
combination of that chart's normals. The absolute coefficients, summed
over all normals and target components, have total budget $Cm(\log m)^2$. Small projections in the appropriate charts contribute
a factor $C\sqrt{(\log m)/m}$, bounding the sum of the Euclidean norms
of the scalar gradients. A dimension-free
first-moment inequality on the cube, proved by the Maurey--Pisier
Gaussian rotation method~\cite{Pis86}, converts this bound into
\[
 \E_{\theta,\eta\in T}
 \|F(v_1(\theta))-F(v_1(\eta))\|_1
 \le C\sqrt m(\log m)^{5/2}.
\]
\Cref{sec:contraction} proves this estimate uniformly in the target
dimension $D$. In particular, the favorable charts may depend on the
scalar component and parameter; the proof averages over charts before
summing the components.

\paragraph{Uniform copies and capacities.}
The probability measure used for these averages is supported on the
first chart, but all charts impose constraints on $F$. We approximate
that measure by multiplicities while retaining at least one copy of
\emph{every} vertex of $V$. Copies of the same vertex have distance zero,
so each contraction on the copies induces a contraction on all of $V$.
The measure on the copies is exactly uniform; its pushforward only
approximates the original first-chart measure. The average bounds survive, and cut-cone duality supplies capacities
with unit demand on each distinct pair of copies.

\Cref{fig:construction} summarizes the two estimates used by this final
step. The size and gap conversion in \Cref{sec:reduction} is
\[
 \mathrm{gap}\gtrsim\frac{m}{\sqrt m(\log m)^{5/2}},
 \qquad \log n\asymp m\log m,
\]
which yields the power $3$ of $\log\log n$ in \Cref{thm:main}.

\begin{figure}[tbp]
\centering
\begin{tikzpicture}[
  >=Latex,
  every node/.style={font=\small,align=center},
  box/.style={draw=blue!45!black,rounded corners=2pt,fill=blue!3,
              inner xsep=5pt,inner ysep=6pt},
  flow/.style={->,semithick,draw=blue!45!black},
  note/.style={font=\footnotesize,fill=white,inner sep=2pt}
]
\node[box,text width=14.25cm] (cube) at (0,0)
  {\textbf{One parameter cube} $U=(-2,2)^m$\\
   Directions for every chart are fixed before the contraction $F$.};
\node[box,text width=4.05cm,minimum height=2.35cm] (chartone) at (-5,-2.0) {};
\node[box,text width=4.05cm,minimum height=2.35cm] (chartmiddle) at (0,-2.0) {};
\node[box,text width=4.05cm,minimum height=2.35cm] (chartlast) at (5,-2.0) {};
\foreach \cx/\chart/\rotation in {-5/1/0,0/s/25,5/S/-20} {
  \node at (\cx,-1.15) {\textbf{Chart $\chart$}};
  \node at (\cx,-1.55) {$\theta\longmapsto v_{\chart}(\theta)$};
  \begin{scope}[shift={(\cx,-2.55)},scale=.95]
    \clip (-.8,-.55) rectangle (.8,.55);
    \begin{scope}[rotate=\rotation]
      \fill[orange!22] (-.24,-.24) rectangle (.24,.24);
      \foreach \offset in {-.72,-.24,.24,.72} {
        \draw[blue!40!black,thin] (\offset,-1.2) -- (\offset,1.2);
        \draw[blue!40!black,thin] (-1.2,\offset) -- (1.2,\offset);
      }
    \end{scope}
    \fill[black] (.06,.02) circle (1.4pt);
    \node[font=\footnotesize,anchor=west,inner sep=1pt,fill=orange!22]
      at (.10,.04) {$\theta$};
  \end{scope}
  \draw[blue!45!black,thin] (\cx-.76,-3.0725) rectangle (\cx+.76,-2.0275);
}
\draw[flow] ([xshift=-5cm]cube.south) -- (chartone.north);
\draw[flow] (cube.south) -- (chartmiddle.north);
\draw[flow] ([xshift=5cm]cube.south) -- (chartlast.north);
\node[box,text width=14.25cm] (metric) at (0,-4.25)
  {\textbf{Hilbert features + triangle repair: negative-type $d$ on all of $V$}\\
   Same-parameter labels: $d(v_s(\theta),v_t(\theta))\le Cr$\\
   $\|x-x'\|_1=\tau$ in one chart: $d((s,x),(s,x'))\le C\tau(m/N)(\log m)^2$};
\draw[flow] (chartone.south) -- ([xshift=-5cm]metric.north);
\draw[flow] (chartmiddle.south) -- (metric.north);
\draw[flow] (chartlast.south) -- ([xshift=5cm]metric.north);
\node[box,text width=6.45cm,minimum height=1.55cm] (large) at (-3.9,-6.35)
  {\textbf{Large metric average}\\
   First-chart law: $v_1(\theta)$, $\theta\sim\mathrm{Unif}(T)$\\
   $\E d(v_1(\theta),v_1(\eta))\ge cm$};
\node[box,text width=6.45cm,minimum height=1.55cm] (small) at (3.9,-6.35)
  {\textbf{Small average for every contraction}\\
   All charts constrain the smoothed gradients\\
   $\E\|F(v_1(\theta))-F(v_1(\eta))\|_1$\\
   $\le C\sqrt m(\log m)^{5/2}$};
\draw[flow] ([xshift=-3.9cm]metric.south) -- (large.north);
\draw[flow] ([xshift=3.9cm]metric.south) -- (small.north);
\node[box,text width=14.25cm] (copies) at (0,-8.4)
  {\textbf{Uniform measure on copies; at least one copy of every vertex}\\
   Copies of a vertex have distance zero. All original chart constraints survive.};
\draw[flow] (large.south) -- ([xshift=-3.9cm]copies.north);
\draw[flow] (small.south) -- ([xshift=3.9cm]copies.north);
\node[box,text width=14.25cm] (duality) at (0,-9.95)
  {\textbf{Cut-cone duality: capacities with unit demand on every distinct pair}\\
   $\OPT(C)\ge1$, \qquad $0<\GL(C)\le C_1(\log m)^{5/2}/\sqrt m$};
\draw[flow] (copies.south) -- (duality.north);
\end{tikzpicture}

\caption{The construction and its two average-distance bounds.
Each shaded region represents the cell of the same parameter $\theta$ in
one rounding partition. Actual cells are determined by all $N$ threshold
families; only two directions are drawn schematically. Drawn lengths do
not represent $d$, and same-parameter labels are $O(r)$ apart.
First-chart averages use independent uniform $\theta,\eta\in T=(-1,1)^m$;
all charts constrain contractions. Arrows indicate construction or
estimate dependencies. The final copies carry unit demands.}
\label{fig:construction}
\end{figure}

\paragraph{Conventions.}
All logarithms are natural. Constants denoted by $c,C>0$ are absolute
and may change from one occurrence to the next. We work with
sufficiently large integers $m$. Euclidean and Hilbert norms are
denoted by $|\cdot|$, and operator norms by $\|\cdot\|_{\op}$.
The notation $A\lesssim B$ means $A\le CB$ for an absolute constant.
All maps into $\ell_1$ used below have a finite, arbitrary target
dimension, and all estimates are independent of that dimension.

\section{Families of directions}
\label{dir:section}

The construction uses many linear images of the same parameter space.  The
directions defining each image must approximate a Gaussian transform, while
most images must have small projections in any prescribed direction.
Sampling Fourier features and controlling them on finite nets is a standard
route to uniform kernel approximation~\cite{RR07}. Here the same samples
must also control derivative profiles and satisfy a majority-chart
guarantee in every direction.

\begin{lemma}\label{lem:directions}
There is an absolute constant $C$ such that, for every sufficiently large
integer $m$, the following holds.  Set $N=m^6$ and $S=m^3$.  There are vectors
$g_{s,i}\in\R^m$, indexed by $1\le s\le S$ and $1\le i\le N$, with these
properties:
\begin{enumerate}
\item For every $s,i$,
\begin{equation}\label{dir:norms}
 \frac{\sqrt m}{2}\le |g_{s,i}|\le 2\sqrt m.
\end{equation}
Within each family $s$, the vectors have pairwise nonparallel directions.
Every eigenvalue of $N^{-1}\sum_{i=1}^N g_{s,i}g_{s,i}^{T}$ belongs to
$[1/2,2]$.
\item For every $s$,
\begin{equation}\label{dir:gaussian-approximation}
 \sup_{|w|\le m^{42}}
 \left|\frac1N\sum_{i=1}^N
 g_{s,i}\sin\langle g_{s,i},w\rangle
       -w e^{-|w|^2/2}\right|
 \le \frac{2}{\sqrt m}.
\end{equation}
\item For every unit vector $v\in\R^m$, at least $S/2$ of the indices $s$
satisfy
\begin{equation}\label{dir:good-charts}
 \max_{1\le i\le N}|\langle v,g_{s,i}\rangle|
 \le C\sqrt{\log m}.
\end{equation}
\end{enumerate}
\end{lemma}

Independent standard Gaussian vectors give such a family. The frequency
net needed in the proof has $\exp(O(m\log m))$ points, whereas the
failure probability for each scalar empirical estimate is at most
$\exp(-cm^4)$. A separate sphere net and independence across charts give
the last assertion simultaneously for every unit direction. The full
probabilistic argument appears in Appendix~\ref{app:directions}.


Fix one such family for the remainder of the proof. In particular, its
every-direction guarantee is available before any contraction is chosen.
Define
\begin{equation}\label{dir:chart-maps}
 u_{s,i}=\frac{g_{s,i}}{\sqrt m},
 \qquad
 J_s:\R^m\longrightarrow\R^N,
 \qquad
 (J_s\theta)_i=\langle u_{s,i},\theta\rangle.
\end{equation}
The eigenvalues of $J_s^TJ_s$ lie in $[N/(2m),2N/m]$, so $J_s$ has full
column rank.  Its left pseudoinverse
$J_s^+=(J_s^TJ_s)^{-1}J_s^T$ satisfies $J_s^+J_s=I_m$, and $J_sJ_s^+$ is the
orthogonal projection onto the range of $J_s$.  We will use
\begin{equation}\label{eq:pseudoinverse}
 \|J_s\|_{\op}\le\sqrt{\frac{2N}{m}},
 \qquad
 \|J_s^+\|_{\op}\le2\sqrt{\frac mN},
 \qquad
 |J_s^+e_i|\le\frac{4m}{N}\quad(1\le i\le N),
\end{equation}
where $e_i$ is the $i$th standard coordinate vector of $\R^N$.  The final
bound follows from $J_s^+e_i=(J_s^TJ_s)^{-1}u_{s,i}$,
$\|(J_s^TJ_s)^{-1}\|_{\op}\le2m/N$, and $|u_{s,i}|\le2$.

\section{A Hilbert kernel with chart extensions}
\label{ker:sec}

Retain the directions, maps, and pseudoinverses from
\Cref{lem:directions}. We construct common Hilbert features on the parameter
space and extend them to the rounded chart labels. Their squared distances
will approximate Euclidean distance. We will control violations of the
triangle inequalities both globally and in proportion to within-chart
coordinate distance. These are the inputs for the metric construction in
\Cref{sec:metric}.

The parameters in this section are
\begin{equation}
\label{ker:parameters}
 a=m^{-40},\qquad b=m^{40},\qquad
 \varepsilon=m^{-200},\qquad r=m^{-8},\qquad \tau=m^{-2000}.
\end{equation}
The numerical exponents only enforce separation of scales. The cutoffs
$a,b$ approximate Euclidean distance; $\varepsilon$ controls the added
Fourier components; $r$ is the tolerance for the angle repair; and $\tau$
is small enough to absorb every polynomial norm bound below. All choices
are fixed powers of $m$, which is essential for the vertex count.

\subsection{The common kernel}

The Gaussian block below produces the macroscopic Euclidean distance
by integrating positive-definite Gaussian kernels over scales, as in the
classical negative-type framework of Schoenberg~\cite{Sch38}.
The additional Fourier blocks have negligible distance contribution but
will supply the coordinate profiles for the chart extensions.

Write $E(t)=(\cos t,\sin t)$, and let $\gamma_m$ be standard Gaussian measure
on $\R^m$.  Define a map $p:\R^m\to\Hh$ into the real Hilbert space
\[
 \Hh=
 L_2([a,b]\times\R^m,d\sigma\,d\gamma_m;\R^2)
 \mathbin{\oplus}
 \bigoplus_{s=1}^{S}\bigoplus_{i=1}^{N}
 L_2([a,b],d\sigma/\sigma;\R^2)
\]
by assigning to $p(\theta)$ the first component
$m^{1/4}E(\langle g,\theta\rangle/\sigma)$ and the component
$\sqrt\varepsilon E(\langle g_{s,i},\theta\rangle/\sigma)$ in summand
$(s,i)$.  The Gaussian characteristic function gives
\begin{equation}
\label{eq:kernel}
\begin{split}
 \langle p(\theta),p(\eta)\rangle
 &=\sqrt m\int_a^b
     e^{-|\eta-\theta|^2/(2\sigma^2)}\,d\sigma\\
 &\quad+\varepsilon\sum_{s=1}^{S}\sum_{i=1}^{N}
   \int_a^b
   \cos\!\left(\frac{\langle g_{s,i},\eta-\theta\rangle}{\sigma}\right)
   \frac{d\sigma}{\sigma}.
\end{split}
\end{equation}

\begin{lemma}[Approximation of Euclidean distance]
\label{lem:kernel}
There is an absolute constant
\[
 c_*=2\int_0^\infty\bigl(1-e^{-1/(2t^2)}\bigr)\,dt>0
\]
such that, whenever $|\theta|,|\eta|\le3\sqrt m$,
\begin{equation}
\label{ker:distance}
 \left|\,|p(\theta)-p(\eta)|^2
       -c_*\sqrt m\,|\theta-\eta|\,\right|\le m^{-37}.
\end{equation}
\end{lemma}

\begin{proof}
The integral defining $c_*$ converges: its integrand is bounded near zero
and is at most $1/(2t^2)$ for $t\ge1$.  Scaling the integration variable
shows that the first component of the squared distance, if integrated over
all positive scales, would equal $c_*\sqrt m\,|\theta-\eta|$.
The omitted scales contribute at most
\[
 2\sqrt m\,a+\frac{\sqrt m\,|\theta-\eta|^2}{b}
 \le 2\sqrt m\,a+\frac{36m^{3/2}}{b}.
\]
The squared distance in all the additional components is at most
$4\varepsilon SN\log(b/a)$.  Consequently the error is at most
\[
 2m^{-79/2}+36m^{-77/2}+320m^{-191}\log m,
\]
which is at most $m^{-37}$ for sufficiently large $m$.
\end{proof}

\subsection{Extensions with small coordinate profiles}

The auxiliary components have amplitude $\sqrt\varepsilon$, while the
preliminary extension below has amplitude $\varepsilon^{-1/2}$. Their
inner products therefore retain full strength despite the negligible
contribution to squared distances. We use this to extend the derivative
of $p$ from each chart subspace with small coordinate profiles against
all the common features. The extension must agree exactly with $Dp$ along
the chart subspace; its small pairings against $p(\eta)$ will control the
triangle inequalities after rounding. We keep separate polynomial bounds
on ordinary Hilbert norms to control the rounding errors.

\begin{lemma}[Chart derivatives]
\label{ker:extension}
For every $s$ there is a smooth map
$A_s:\R^m\to\mathcal L(\R^N,\Hh)$ such that
\begin{equation}
\label{ker:extension-identity}
 A_s(\theta)J_s=Dp(\theta).
\end{equation}
For $|\theta|,|\eta|\le3\sqrt m$ and $1\le i\le N$,
\begin{equation}
\label{ker:profile}
 \bigl|\langle A_s(\theta)e_i,p(\eta)\rangle\bigr|
 \le C\frac{m}{N}\log m.
\end{equation}
Moreover, uniformly in $\theta\in\R^m$ and $s$,
\begin{equation}
\label{ker:strong-bounds}
 |p(\theta)|,\quad \|Dp(\theta)\|_{\op},\quad
 \|D^2p(\theta)\|_{\op},\quad
 \|A_s(\theta)\|_{\op},\quad
 \|DA_s(\theta)\|_{\op}\le m^{300}.
\end{equation}
Here the norm of a second derivative or an operator-valued derivative is
the corresponding multilinear operator norm.
\end{lemma}

\begin{proof}
We first identify the profile that the preliminary extension should
reproduce. For $|\theta|,|\eta|\le3\sqrt m$, put
$w=(\eta-\theta)/\sigma$. Differentiation in $\theta$ of the Gaussian
term in \eqref{eq:kernel} gives the vector
\[
 \sqrt m\int_a^b w e^{-|w|^2/2}\,\frac{d\sigma}{\sigma}.
\]
The empirical approximation in \Cref{lem:directions} suggests replacing
$w e^{-|w|^2/2}$ by $N^{-1}\sum_i g_{s,i}\sin\langle g_{s,i},w\rangle$.
We realize this profile using the auxiliary Fourier components.
Define $A_s^0(\theta)e_i$ to be supported in summand $(s,i)$ of $\Hh$,
where it is the function
\[
 \frac{m}{N}\varepsilon^{-1/2}
 E'\!\left(\frac{\langle g_{s,i},\theta\rangle}{\sigma}\right).
\]
Its profile is
\begin{equation}
\label{ker:initial-profile}
 \langle A_s^0(\theta)e_i,p(\eta)\rangle
 =\frac{m}{N}\int_a^b\sin\langle g_{s,i},w\rangle\,
          \frac{d\sigma}{\sigma}.
\end{equation}
Since $u_{s,i}=g_{s,i}/\sqrt m$, the profile of
$A_s^0J_s$ is represented by the vector
\[
 \sqrt m\int_a^b
 \left(\frac1N\sum_{i=1}^N g_{s,i}\sin\langle g_{s,i},w\rangle\right)
 \frac{d\sigma}{\sigma}.
\]
Thus $A_s^0J_s$ has the intended empirical profile. Set
$R_s=Dp-A_s^0J_s$; we now bound its pairings against the common features
before correcting the extension on the chart subspace.
For the parameters under consideration,
$|w|\le6\sqrt m/a<m^{42}$.  Thus \Cref{lem:directions} bounds the
difference between these vectors by $2\log(b/a)$.
The gradient of the second term of \eqref{eq:kernel} has norm at most
\[
 \varepsilon\sum_{t=1}^{S}\sum_{i=1}^{N}|g_{t,i}|
        \int_a^b\frac{d\sigma}{\sigma^2}
 \le\frac{2\varepsilon SN\sqrt m}{a}\le1.
\]
It follows that
\begin{equation}
\label{ker:residual-profile}
 \bigl|\langle R_s(\theta)v,p(\eta)\rangle\bigr|
 \le C(\log m)|v|.
\end{equation}
Now define
\begin{equation}
\label{ker:A-definition}
 A_s=A_s^0+R_sJ_s^+=A_s^0(I-J_sJ_s^+)+Dp\,J_s^+.
\end{equation}
The identity $J_s^+J_s=I$ proves \eqref{ker:extension-identity}.
Combining \eqref{ker:initial-profile}, \eqref{ker:residual-profile}, and
the column bound $|J_s^+e_i|\le4m/N$ from \eqref{eq:pseudoinverse}
proves \eqref{ker:profile}.

For completeness, the required Hilbert norm bounds have ample polynomial
slack.  Differentiation under the Hilbert space formulas is valid to every
order, because $a>0$ and Gaussian moments of every order are finite.
For $j=0,1,2$, the squared multilinear derivative norm of the first
component of $p$ is at most
\[
 \sqrt m\,b\,a^{-2j}\E|G|^{2j},
\]
and that of all the additional components together is at most
\[
 \varepsilon SN\log(b/a)
       \left(\frac{2\sqrt m}{a}\right)^{2j}.
\]
Using $\E|G|^2=m$ and $\E|G|^4=m(m+2)$, these estimates give
$\|D^jp\|_{\op}\le m^{110}$ for $0\le j\le2$.
Bounding an operator by $\sqrt N$ times its largest column norm gives
\[
 \|A_s^0\|_{\op}
 \le \sqrt N\,\frac{m}{N}\varepsilon^{-1/2}\sqrt{\log(b/a)},
 \qquad
 \|DA_s^0\|_{\op}
 \le \sqrt N\,\frac{m}{N}\varepsilon^{-1/2}
          \sqrt{\log(b/a)}\frac{2\sqrt m}{a}.
\]
In particular, these two quantities are at most $m^{110}$ and $m^{150}$,
respectively.  Substitute these bounds in
\eqref{ker:A-definition}, together with
$R_s=Dp-A_s^0J_s$, its derivative, and
$\|J_s\|_{\op}\le2\sqrt N$ and
$\|J_s^+\|_{\op}\le2\sqrt{m/N}$.
The projection $I-J_sJ_s^+$ has norm at most one; the last expression
in \eqref{ker:A-definition} and its derivative therefore bound all five
quantities in \eqref{ker:strong-bounds} by $m^{300}$ for sufficiently
large $m$.
\end{proof}

\subsection{Rounded charts and their angles}

Let $Q_s=I-J_sJ_s^+$, the orthogonal projection perpendicular to the
image of $J_s$.  Define
\begin{equation}
\label{ker:P-definition}
 P_s(x)=p(J_s^+x)+A_s(J_s^+x)Q_sx,
 \qquad x\in\R^N.
\end{equation}
On the chart subspace this gives $P_s(J_s\theta)=p(\theta)$ and
$DP_s(J_s\theta)=A_s(\theta)$: use $Q_sJ_s=0$ and
\eqref{ker:extension-identity}. We use only values very close to that
subspace; the derivative away from it is computed explicitly below.
Let $U=(-2,2)^m$, and for $\theta\in U$ put
\begin{equation}
\label{ker:rounding}
 X_s(\theta)_i
 =\tau\left\lfloor\frac{\langle u_{s,i},\theta\rangle}{\tau}
       \right\rfloor.
\end{equation}
In each chart, discard the threshold hyperplanes
$\langle u_{s,i},\theta\rangle=k\tau$, $k\in\mathbb Z$;
their union has Lebesgue measure zero.  Define the finite set
\begin{equation}
\label{ker:vertices}
 V=\{(s,X_s(\theta)):1\le s\le S,\ 
       \theta\in U\text{ off the threshold hyperplanes of chart }s\}.
\end{equation}
Write $v_s(\theta)=(s,X_s(\theta))$ whenever this is defined, and
$P_v=P_s(x)$ for $v=(s,x)\in V$.

For any Hilbert vectors $Z_v$, their squared distances satisfy all
triangle inequalities exactly when
\begin{equation}\label{rep:eq:angle}
\langle Z_v-Z_z,Z_y-Z_z\rangle\ge0
\qquad(v,y,z\in V).
\end{equation}
We next bound how far the vectors $P_v$ can violate this condition.

\begin{lemma}[Two angle estimates]
\label{ker:angles}
For sufficiently large $m$, all $v,z,y\in V$ satisfy
\begin{equation}
\label{ker:global-angle}
 \langle P_v-P_z,P_y-P_z\rangle\ge-r.
\end{equation}
There is an absolute constant $C_0$ such that, with $L=C_0\log m$, for
$v=(s,x)$ and $z=(s,x')$ in the same chart and every $y\in V$,
\begin{equation}
\label{ker:local-angle}
 \bigl|\langle P_v-P_z,P_y-P_z\rangle\bigr|
 \le L\frac{m}{N}\|x-x'\|_1.
\end{equation}
Moreover, for all defined $v_s(\theta),v_t(\eta)$,
\begin{equation}
\label{ker:rounded-distance}
 \left|\,|P_{v_s(\theta)}-P_{v_t(\eta)}|^2
       -c_*\sqrt m\,|\theta-\eta|\,\right|
 \le m^{-37}+m^{610}\tau\le r/10.
\end{equation}
\end{lemma}

\begin{proof}
For $x=X_s(\theta)$, rounding and \eqref{eq:pseudoinverse} give
\begin{equation}
\label{ker:rounding-errors}
 |x-J_s\theta|\le\sqrt N\tau,\qquad
 |J_s^+x-\theta|\le2\sqrt m\tau,\qquad
 |Q_sx|\le\sqrt N\tau.
\end{equation}
In particular, $|J_s^+x|\le3\sqrt m$.  By
\eqref{ker:strong-bounds} and \eqref{ker:P-definition},
\begin{equation}
\label{ker:feature-error}
 |P_{v_s(\theta)}-p(\theta)|
 \le m^{300}(2\sqrt m+\sqrt N)\tau
 \le m^{305}\tau.
\end{equation}
The change in a squared distance caused by the two errors in
\eqref{ker:feature-error} is at most $m^{610}\tau$, since
$|p(\theta)|\le m^{300}$.  Therefore \Cref{lem:kernel} gives
\eqref{ker:rounded-distance}; its last inequality follows directly from
\eqref{ker:parameters}.

Choose parameter representatives in $U$ for $v,z,y$.  Apply
\eqref{ker:rounded-distance} to their three pairs and use
\[
 2\langle P_v-P_z,P_y-P_z\rangle
 =|P_v-P_z|^2+|P_y-P_z|^2-|P_v-P_y|^2.
\]
The corresponding expression for the ordinary Euclidean distances of
the three representatives is nonnegative by the triangle inequality.
The total error is at most $3r/10$, proving
\eqref{ker:global-angle}.

For the local estimate, we integrate coordinate profiles so that the bound
scales with $\|x-x'\|_1$. Consider the segment joining $x$ to $x'$.
Every point $\xi$ on this segment satisfies
$|Q_s\xi|\le\sqrt N\tau$ and $|J_s^+\xi|\le3\sqrt m$, by convexity
and \eqref{ker:rounding-errors}.  Differentiating
\eqref{ker:P-definition} and using
\eqref{ker:extension-identity} gives the exact identity
\begin{equation}
\label{ker:DP-identity}
 DP_s(\xi)h
 =A_s(J_s^+\xi)h
  +DA_s(J_s^+\xi)[J_s^+h]Q_s\xi.
\end{equation}
The operator norm of the second term is at most
$2m^{300}\sqrt m\tau\le m^{305}\tau$.
For a vertex $y=v_t(\eta)$, replace $P_y$ by $p(\eta)$ using
\eqref{ker:feature-error}.  Equations \eqref{ker:profile} and
\eqref{ker:strong-bounds} then imply, throughout this segment,
\[
 \bigl|\langle DP_s(\xi)e_i,P_y\rangle\bigr|
 \le C\frac{m}{N}\log m+3m^{605}\tau
 \le C'\frac{m}{N}\log m.
\]
The same bound holds with $P_z$ in place of $P_y$.
Integrating \eqref{ker:DP-identity} along the segment and summing the
absolute coordinate increments proves \eqref{ker:local-angle} after
choosing a sufficiently large absolute $C_0$.
\end{proof}

\section{The finite negative-type semimetric}
\label{sec:repair}
\label{sec:metric}

The vectors $P_v$ from \Cref{ker:angles} have the desired macroscopic
distances. Their inner products in \eqref{rep:eq:angle} can still be
negative, but the lemma bounds their negative parts both by $r$ and
by a quantity proportional to within-chart coordinate distance.
We complete the construction by appending a Hilbert correction that
compensates the smaller of these bounds. A correction controlled only
by $r$ could dominate the cost of crossing an individual cell interface.

The following lemma corrects violations of \eqref{rep:eq:angle} while
preserving small distances up to a logarithmic loss. We state it for
arbitrary chartwise $\ell_1$ distances so that the required local and
global bounds are explicit.

\begin{lemma}[Triangle repair]\label{lem:repair}
Let a finite set $V$ be partitioned into charts. In each chart let
$\delta$ be a nonnegative multiple of an $\ell_1$ distance, and set
$\delta(v,w)=\infty$ for points in different charts.
Suppose that $r>0$ and Hilbert vectors $P_v$ satisfy
\begin{equation}\label{rep:eq:hypothesis}
\langle P_v-P_z,P_y-P_z\rangle
\ge-\min\{r,\delta(v,z),\delta(y,z)\}
\qquad(v,y,z\in V).
\end{equation}
Then there are Hilbert vectors $q_v$ such that
\[
d(v,w)=|P_v-P_w|^2+|q_v-q_w|^2
\]
is a negative-type semimetric. They satisfy
\begin{equation}\label{rep:eq:distance}
|q_v-q_w|^2
=8\int_0^r\bigl(1-e^{-\delta(v,w)/l}\bigr)\,dl\le8r.
\end{equation}
If $0<\delta(v,w)\le r$, then
\begin{equation}\label{rep:eq:local}
|q_v-q_w|^2
\le8\delta(v,w)\left(1+\log\frac r{\delta(v,w)}\right).
\end{equation}
If $\delta(v,w)=0$, then $q_v=q_w$.
\end{lemma}

\begin{proof}
We use the convention $e^{-\infty}=0$. Define a Gram matrix by
\begin{equation}\label{rep:eq:gram}
\langle q_v,q_w\rangle=4\int_0^r e^{-\delta(v,w)/l}\,dl.
\end{equation}
This matrix is positive semidefinite. Indeed, in one coordinate the
map $t\mapsto\mathbf1_{[b,t]}$ into $L_2$ has squared distance
$|t-t'|$, where $b$ is a fixed lower bound on the finitely many
coordinate values. Direct sums and scaling show that an $\ell_1$
distance is squared Hilbertian. For finitely many Hilbert vectors
$h_v$, the kernel
\[
e^{-|h_v-h_w|^2/l}
=\E\cos\!\left(\sqrt{2/l}\,\langle G,h_v-h_w\rangle\right)
\]
is positive semidefinite: $G$ is a standard Gaussian in their finite
linear span, and the cosine is the inner product of the corresponding
$(\cos,\sin)$ pairs. This is the Gaussian-kernel criterion of
Schoenberg~\cite{Sch38}, here in its elementary finite form.
The kernels in \eqref{rep:eq:gram} are block diagonal across charts;
integration preserves positive semidefiniteness. Thus the specified
Gram matrix has a Hilbert realization.

The extended distance $\delta$ satisfies the triangle inequality,
including triples meeting different charts. Consequently, for
$a=\delta(v,z)$ and $b=\delta(y,z)$,
\[
\begin{split}
1-e^{-a/l}-e^{-b/l}+e^{-\delta(v,y)/l}
&\ge 1-e^{-a/l}-e^{-b/l}+e^{-(a+b)/l}\\
&=(1-e^{-a/l})(1-e^{-b/l}).
\end{split}
\]
For $t=\min\{r,a,b\}$, this product is at least $(1-e^{-1})^2$
when $0<l<t$. Equation \eqref{rep:eq:gram} therefore gives
\[
\langle q_v-q_z,q_y-q_z\rangle
\ge4(1-e^{-1})^2t\ge t.
\]
Together with \eqref{rep:eq:hypothesis}, this proves
\eqref{rep:eq:angle} for $Z_v=(P_v,q_v)$.

The diagonal entries of \eqref{rep:eq:gram} are $4r$, which proves
\eqref{rep:eq:distance}. Finally, $1-e^{-s}\le\min\{1,s\}$ and
\[
\int_0^r\min\{1,t/l\}\,dl=t\bigl(1+\log(r/t)\bigr)
\qquad(0<t\le r)
\]
give \eqref{rep:eq:local}. The assertion at distance zero follows
directly from \eqref{rep:eq:distance}.
\end{proof}

We now apply \Cref{lem:repair} to the rounded labels of \Cref{ker:angles}.
The following proposition collects the metric properties needed for
the contraction and duality arguments.

\begin{proposition}\label{prop:metric}
For every sufficiently large integer $m$, the vertex set $V$ in
\eqref{ker:vertices} admits a negative-type semimetric $d$ with the
following properties. Its chart maps $v_s(\theta)$, defined almost
everywhere on $U=(-2,2)^m$, satisfy
\begin{equation}\label{met:eq:macro}
\left|d(v_s(\theta),v_t(\eta))
-c_*\sqrt m\,|\theta-\eta|\right|\le Cr
\qquad(\theta,\eta\in U)
\end{equation}
whenever the labels are defined, where $r=m^{-8}$.
For labels in the same chart,
\begin{equation}\label{met:eq:adjacent}
\|x-x'\|_1=\tau
\quad\Longrightarrow\quad
d((s,x),(s,x'))\le C\tau\frac mN(\log m)^2.
\end{equation}
Moreover,
\begin{equation}\label{met:eq:size-diameter}
|V|\le\exp(Cm\log m),\qquad
\operatorname{diam}(V,d)\le Cm,
\end{equation}
and, for independent uniform $\theta,\eta\in T=(-1,1)^m$,
\begin{equation}\label{met:eq:average}
\E d(v_1(\theta),v_1(\eta))\ge cm.
\end{equation}
Here $N=m^6$, $\tau=m^{-2000}$, and all constants are absolute.
\end{proposition}

\begin{proof}
Let $P_v$ and $L=C_0\log m$ be as in \Cref{ker:angles}. Within a
chart set
\[
\delta((s,x),(s,x'))=L\frac mN\|x-x'\|_1,
\]
and put $\delta=\infty$ between different charts.
Equation \eqref{ker:global-angle} bounds every $P$-angle below by
$-r$. Equation \eqref{ker:local-angle} also bounds it below by
$-\delta(v,z)$ whenever $v,z$ are in the same chart, and by
$-\delta(y,z)$ whenever $y,z$ are in the same chart, after swapping
the two arguments of the inner product. Thus the hypothesis of
\Cref{lem:repair} holds. Use its vectors $q_v$ and set
\begin{equation}\label{met:eq:definition}
d(v,w)=|(P_v,q_v)-(P_w,q_w)|^2.
\end{equation}
This is a negative-type semimetric. The bound
$|q_v-q_w|^2\le8r$ and \eqref{ker:rounded-distance} prove
\eqref{met:eq:macro}.

Taking $y=v$ in \eqref{ker:local-angle} gives
$|P_v-P_z|^2\le\delta(v,z)$. If $\|x-x'\|_1=\tau$, then
\[
\delta= C_0\tau(m/N)\log m<r,
\qquad
\log(r/\delta)
=1997\log m-\log(C_0\log m)=O(\log m).
\]
Equation \eqref{rep:eq:local} now proves
\eqref{met:eq:adjacent}.

We next count the rounded labels. Since $|u_{s,i}|\le2$, a
hyperplane $\langle u_{s,i},\theta\rangle=k\tau$ meeting $U$ has
$|k\tau|<4\sqrt m$. Hence at most
\[
H\le N(1+8\sqrt m/\tau)
\]
threshold hyperplanes meet $U$ in a chart. An arrangement of $H$
hyperplanes in $\R^m$ has at most
$\sum_{j=0}^m\binom Hj$ full-dimensional regions (see~\cite{Zas75}).
This follows by
induction: an added hyperplane creates at most as many new regions
as the preceding hyperplanes cut out on it. The labels are constant
on the regions, and their intersections with the convex cube $U$
are connected. Therefore
\[
|V|\le S\sum_{j=0}^m\binom Hj\le\exp(Cm\log m),
\]
because $H$ is bounded by a fixed power of $m$. The diameter bound
in \eqref{met:eq:size-diameter} follows from
\eqref{met:eq:macro} and $\operatorname{diam}(U)=4\sqrt m$.

Finally, for independent uniform parameters on $T$,
\[
\E|\theta-\eta|^2=\frac{2m}{3},\qquad
|\theta-\eta|\le2\sqrt m.
\]
It follows that $\E|\theta-\eta|\ge\sqrt m/3$.
Equation \eqref{met:eq:macro} gives
$\E d(v_1(\theta),v_1(\eta))\ge c_*m/3-Cr$, proving
\eqref{met:eq:average} for sufficiently large $m$.
\end{proof}

The macroscopic comparison \eqref{met:eq:macro} does not make the
interface constraints redundant: its additive error $O(r)=O(m^{-8})$
is far larger than the interface bound $O(m^{-2005}(\log m)^2)$ in
\eqref{met:eq:adjacent}. The next section uses these much smaller
interface distances to constrain every contraction.

\section{Contractions into \texorpdfstring{$\ell_1$}{ell1}}
\label{sec:contraction}

The small jumps in every chart force small average variation in $\ell_1$.
Throughout this section, scalar gradients carry the Euclidean norm.

\begin{proposition}[Contraction bound]
\label{prop:contraction}
Let $(V,d)$ and the chart maps $v_s:U\to V$ be those of
\Cref{prop:metric}, where $U=(-2,2)^m$, $1\le s\le S=m^3$, and
$N=m^6$. There is an absolute constant $C$ such that every map
$F:V\to\ell_1^D$, for every finite $D$, satisfying
\[
  \|F(v)-F(w)\|_1\le d(v,w)\qquad(v,w\in V)
\]
obeys
\begin{equation}
  \E_{\theta,\eta\in T}
  \|F(v_1(\theta))-F(v_1(\eta))\|_1
  \le C\sqrt m\,(\log m)^{5/2},
  \label{con:eq:contraction}
\end{equation}
where $T=(-1,1)^m$ and the two parameters are independent and uniform.
\end{proposition}

After smoothing, each gradient is a combination of chart normals. The
jump bounds give a total coefficient budget $Cm(\log m)^2$, while the
good-chart property contributes the factor $C\sqrt{(\log m)/m}$.
We then use the following scalar inequality to pass from gradients to
average distances without a further dimension factor.

\begin{lemma}[First-moment inequality on a cube]
\label{con:lem:cube}
Let $T=(-1,1)^m$, and let $h$ be a real-valued $C^1$ function on a
neighborhood of $\overline T$. If $\theta,\eta$ are independent and uniform
on $T$, then
\begin{equation}
  \E|h(\theta)-h(\eta)|
  \le \E|\nabla h(\theta)|.
  \label{con:eq:cube}
\end{equation}
\end{lemma}

We use the Gaussian rotation argument of Maurey and Pisier~\cite{Pis86},
followed by the coordinatewise Gaussian distribution function.

\begin{proof}
Write $\Phi$ for the standard Gaussian distribution function, set
$\psi(t)=2\Phi(t)-1$, and let $\Psi:\R^m\to T$ apply $\psi$
coordinatewise. If $G,G'$ are independent standard Gaussian vectors in
$\R^m$, then $\Psi(G),\Psi(G')$ are independent and uniform on $T$.
For $0\le t\le\pi/2$, put
\[
  Y(t)=\cos(t)G+\sin(t)G',
  \qquad
  Y'(t)=-\sin(t)G+\cos(t)G'.
\]
At each fixed time, $Y(t)$ and $Y'(t)$ are independent standard Gaussian
vectors. Consequently, conditioning on $Y(t)$ gives
\begin{align*}
  \E\left[\left|\frac{d}{dt}h(\Psi(Y(t)))\right|
      \,\middle|\,Y(t)\right]
  &=\sqrt{\frac{2}{\pi}}\,
    \left|D\Psi(Y(t))^{T}\nabla h(\Psi(Y(t)))\right|\\
  &\le \frac{2}{\pi}\,|\nabla h(\Psi(Y(t)))|,
\end{align*}
because $\sup|\psi'|=\sqrt{2/\pi}$. The fundamental theorem of calculus,
followed by Tonelli's theorem, bounds the expected difference between the
endpoints by the integral of the expected absolute derivative. The
distribution of $\Psi(Y(t))$ is uniform on $T$ for every $t$, so integration
over $[0,\pi/2]$ proves \eqref{con:eq:cube}. All the integrands are
integrable: $\nabla h$ is bounded on $\overline T$, and Gaussian vectors have
finite first moments.
\end{proof}

\begin{proof}[Proof of \Cref{prop:contraction}]
We use $r=m^{-8}$ and $\tau=m^{-2000}$ as in \Cref{prop:metric}.
For each chart, define
\[
  f_s(\theta)=F(v_s(\theta))\quad\text{for almost every }\theta\in U,
\]
and extend $f_s$ by zero outside $U$. Set $f_s=0$ also on the threshold hyperplanes and on $\partial U$.
These are bounded, compactly supported functions because $V$ is finite. The
macroscopic estimate in \Cref{prop:metric} implies
\begin{align}
  \|f_s(\theta)-f_1(\theta)\|_1&\le Cr,
       &&\text{for almost every }\theta\in U,\label{con:eq:chart-close}\\
  \|f_s(\theta)-f_s(\eta)\|_1
       &\le c_*\sqrt m\,|\theta-\eta|+Cr,
       &&\text{off the threshold hyperplanes in }U.
       \label{con:eq:macro-variation}
\end{align}

\paragraph{Smoothing.}
Fix an even, nonnegative function $\rho\in C_c^\infty((-1,1))$ with
$\int_\R\rho=1$. Set $\lambda=m^{-2}$ and
\[
  \kappa_\lambda(z)=\prod_{j=1}^m\lambda^{-1}\rho(z_j/\lambda),
  \qquad H_s=f_s*\kappa_\lambda.
\]
The functions $H_s:\R^m\to\R^D$ are smooth. For sufficiently large $m$,
the support of $x\mapsto\kappa_\lambda(\theta-x)$ lies in $U$ whenever
$\theta\in\overline T$. Also
\[
  \int_{\R^m}|\partial_j\kappa_\lambda|
       =\lambda^{-1}\|\rho'\|_{L_1(\R)}.
\]
Every displacement in the support of $\kappa_\lambda$ has Euclidean norm
at most $\sqrt m\lambda$. Thus \eqref{con:eq:macro-variation} gives
\begin{equation}
  \|H_s(\theta)-f_s(\theta)\|_1\le C(m\lambda+r)
  \quad\text{for almost every }\theta\in T.
  \label{con:eq:smoothing-error}
\end{equation}
Since $f_s-f_1$ vanishes outside $U$, \eqref{con:eq:chart-close} bounds its
$\ell_1$ norm by $Cr$ almost everywhere on $\R^m$. Differentiating the
convolution therefore yields, for every $\theta\in T$ and every $j$,
\begin{equation}
  \|\partial_jH_s(\theta)-\partial_jH_1(\theta)\|_1
      \le Cr/\lambda.
  \label{con:eq:derivative-mismatch}
\end{equation}

\paragraph{The derivative carried by each family of interfaces.}
Fix a chart $s$ and abbreviate $u_i=u_{s,i}$. Its interfaces are contained
in the hyperplanes
\[
  Q_{i,k}=\{x\in\R^m:\langle u_i,x\rangle=k\tau\},
  \qquad 1\le i\le N,\quad k\in\mathbb Z.
\]
Only finitely many of these hyperplanes meet $U$. At almost every point
of $Q_{i,k}\cap U$, let $\Delta_{i,k}(x)\in\R^D$ denote the value of
$f_s$ on the side of increasing $\langle u_i,x\rangle$ minus its value
on the opposite side. This is well-defined away from intersections with
the other interface hyperplanes. By the nonparallelism in
\Cref{lem:directions}, those intersections have zero
$(m-1)$-dimensional measure. The two labels at such an interface differ
by $\tau e_i$, so the adjacent-label estimate in \Cref{prop:metric} gives
\begin{equation}
  \|\Delta_{i,k}(x)\|_1
      \le C\tau\frac{m}{N}(\log m)^2
  \quad\text{for almost every }x\in Q_{i,k}\cap U.
  \label{con:eq:jump-bound}
\end{equation}

For completeness, the distributional gradient of each scalar component
of $f_s$ on $U$ is the vector-valued measure
\begin{equation}
  D f_{s,\alpha}
   =\sum_{i,k}\frac{u_i}{|u_i|}\,
       \Delta_{i,k,\alpha}\,
       \mathcal H^{m-1}\!\restriction(Q_{i,k}\cap U).
  \label{con:eq:distributional-gradient}
\end{equation}
Indeed, test against a smooth function compactly supported in $U$ and
integrate by parts on each polyhedral cell. The cell interiors contribute
nothing because $f_s$ is constant there. The two contributions on each
shared facet combine into the jump times the unit normal $u_i/|u_i|$;
lower-dimensional intersections have zero facet measure. The test
function vanishes near $\partial U$, so there is no contribution from
the outer boundary.

For $\theta\in T$, the kernel $\kappa_\lambda(\theta-\cdot)$ is a smooth
test function compactly supported in $U$. Convolving
\eqref{con:eq:distributional-gradient} with this kernel gives, for every
component $\alpha$,
\begin{equation}
  \nabla H_{s,\alpha}(\theta)
       =\sum_{i=1}^N u_{s,i}\,b_{s,i,\alpha}(\theta),
  \label{con:eq:gradient-representation}
\end{equation}
where the vector coefficient $b_{s,i}(\theta)\in\R^D$ is
\begin{equation}
  b_{s,i}(\theta)=\frac{1}{|u_{s,i}|}
       \sum_{k\in\mathbb Z}\int_{Q_{i,k}\cap U}
       \kappa_\lambda(\theta-x)\Delta_{i,k}(x)
       \,d\mathcal H^{m-1}(x).
  \label{con:eq:coefficient-formula}
\end{equation}
This identity holds at every $\theta\in T$: distributional convolution
agrees with the classical derivative of the smooth function $H_s$.

We claim that the coefficients satisfy the pointwise budget
\begin{equation}
  \sum_{\alpha=1}^D|b_{s,i,\alpha}(\theta)|
       \le C\frac{m}{N}(\log m)^2
       \qquad(\theta\in T).
  \label{con:eq:coefficient-budget}
\end{equation}
To establish it, fix $u=u_{s,i}$ and $\theta$, and let $Z$ have density
$x\mapsto\kappa_\lambda(\theta-x)$. Write $h$ for the density of the
real random variable $\langle u,Z\rangle$. The coarea formula for this
linear functional states that
\[
  h(t)=\frac{1}{|u|}
       \int_{\langle u,x\rangle=t}
       \kappa_\lambda(\theta-x)\,d\mathcal H^{m-1}(x).
\]
The lower bound $|u|\ge1/2$ in \Cref{lem:directions} supplies a coordinate
$j$ with $|u_j|\ge1/(2\sqrt m)$. Conditional on the other coordinates of
$Z$, the variable $\langle u,Z\rangle$ has a translated, possibly
reflected copy of the one-dimensional density $\rho$, scaled by
$\lambda|u_j|$. Its derivative has $L_1$ norm
$\|\rho'\|_{L_1}/(\lambda|u_j|)$. Taking the mixture over the other
coordinates shows that $h$ is continuously differentiable, compactly
supported, and satisfies
\begin{equation}
  \int_\R|h'(t)|\,dt\le C\sqrt m/\lambda.
  \label{con:eq:density-variation}
\end{equation}

For any nonnegative, continuously differentiable, compactly supported
density $h$, comparison on each interval $[k\tau,(k+1)\tau]$ gives
\[
  \tau h(k\tau)
     \le\int_{k\tau}^{(k+1)\tau}h(t)\,dt
       +\tau\int_{k\tau}^{(k+1)\tau}|h'(t)|\,dt.
\]
Summing and using \eqref{con:eq:density-variation}, we obtain
\begin{equation}
  \tau\sum_{k\in\mathbb Z}h(k\tau)
       \le1+C\tau\sqrt m/\lambda\le2
  \label{con:eq:lattice-bound}
\end{equation}
for all sufficiently large $m$. Because the convolution kernel is
nonnegative and is supported inside $U$, the triangle inequality in
\eqref{con:eq:coefficient-formula}, followed by
\eqref{con:eq:jump-bound} and \eqref{con:eq:lattice-bound}, proves
\eqref{con:eq:coefficient-budget}.

\paragraph{Combining the charts.}
Fix $\theta\in T$ and a component $\alpha$, and put
$a_\alpha=\nabla H_{1,\alpha}(\theta)$. If $a_\alpha\ne0$, the
direction property in \Cref{lem:directions} gives at least $S/2$ charts
$s$ such that
\[
  \max_i\left|\left\langle
       \frac{a_\alpha}{|a_\alpha|},u_{s,i}\right\rangle\right|
       \le C\sqrt{\frac{\log m}{m}}.
\]
For each such chart, \eqref{con:eq:gradient-representation} implies
\[
  |a_\alpha|
    \le C\sqrt{\frac{\log m}{m}}\sum_i|b_{s,i,\alpha}(\theta)|
       +|\nabla H_{s,\alpha}(\theta)-a_\alpha|.
\]
All terms on the right are nonnegative. Summing over the good charts and
then enlarging that sum to all charts therefore gives
\begin{equation}
  |a_\alpha|\le\frac{2}{S}\sum_{s=1}^S
       \left(
       C\sqrt{\frac{\log m}{m}}\sum_i|b_{s,i,\alpha}(\theta)|
       +|\nabla H_{s,\alpha}(\theta)-a_\alpha|
       \right).
  \label{con:eq:good-chart-average}
\end{equation}
For $a_\alpha=0$ the same inequality holds trivially. The good charts may
depend on $\alpha$ and $\theta$; \eqref{con:eq:good-chart-average} removes
that dependence before we sum over components.

For each $s$, the Euclidean norm is at most the sum of absolute
coordinates, so \eqref{con:eq:derivative-mismatch} gives
\[
  \sum_{\alpha=1}^D
       |\nabla H_{s,\alpha}(\theta)-\nabla H_{1,\alpha}(\theta)|
  \le\sum_{j=1}^m
       \|\partial_jH_s(\theta)-\partial_jH_1(\theta)\|_1
  \le Cmr/\lambda.
\]
Sum \eqref{con:eq:good-chart-average} over $\alpha$ and apply this bound
and \eqref{con:eq:coefficient-budget}. We obtain at every $\theta\in T$
\begin{align}
  \sum_{\alpha=1}^D|\nabla H_{1,\alpha}(\theta)|
   &\le C\sqrt{\frac{\log m}{m}}\,
          N\frac{m}{N}(\log m)^2+Cmr/\lambda\notag\\
   &\le C\sqrt m\,(\log m)^{5/2},
   \label{con:eq:total-gradient}
\end{align}
where $mr/\lambda=m^{-5}$.

Finally, apply \Cref{con:lem:cube} to every component of $H_1$ and sum.
Equation \eqref{con:eq:total-gradient} yields
\[
  \E_{\theta,\eta\in T}\|H_1(\theta)-H_1(\eta)\|_1
       \le C\sqrt m\,(\log m)^{5/2}.
\]
The expected distance for $f_1$ exceeds this by at most
$2C(m\lambda+r)$, by \eqref{con:eq:smoothing-error}. Since
$m\lambda=m^{-1}$, this additional term is absorbed into the same bound.
The definition of $f_1$ now proves \eqref{con:eq:contraction}.
\end{proof}

\section{Exactly uniform demands}
\label{sec:reduction}

We finish by converting the metric obstruction into capacities with unit
demand on every unordered pair.  We first record the precise duality
statement, allowing distinct points to have distance zero. This is the
average-distortion form of the cut/embedding correspondence
\cite[Section~4]{LLR95}\cite[Section~1.3]{Rab08}; the explicit programs
below keep the unit-demand normalization visible.

\begin{lemma}[Cut-cone duality]
\label{lem:duality}
Let $d$ be a semimetric on $[n]$, where $n\geq 2$, that satisfies the
triangle inequalities and is a squared Hilbert distance.  Suppose that
$a,b>0$ satisfy
\begin{equation}
\label{red:averages}
\frac{1}{n^2}\sum_{i,j=1}^n d(i,j)\geq a,
\qquad
\frac{1}{n^2}\sum_{i,j=1}^n
\|F(i)-F(j)\|_1\leq b
\end{equation}
for every map $F$ into a finite-dimensional $\ell_1$ space such that
$\|F(i)-F(j)\|_1\leq d(i,j)$ for all $i,j$.
Then there are nonnegative capacities $c_{ij}$ on unordered pairs of
distinct vertices such that, with demand exactly one on every pair,
\begin{equation}
\label{red:gap-lemma}
\OPT(C)\geq 1,
\qquad
0<\GL(C)\leq \frac{b}{a}.
\end{equation}
\end{lemma}

\begin{proof}
Let $\mathcal S$ be the finite collection of nonempty proper subsets of
$[n]$.  For $S\in\mathcal S$, write
$\delta_S(i,j)=|\mathbf 1_S(i)-\mathbf 1_S(j)|$.  Consider the finite
linear program.\footnote{Matou\v{s}ek and
Rabinovich~\cite{MR01} study the related convex hull of individually
dominated line metrics. Here the cut coefficients are arbitrary
nonnegative numbers; only their sum as a distance function is required
to be dominated by $d$.}
\begin{equation}
\label{red:primal}
\begin{aligned}
\text{maximize}\quad &
\sum_{S\in\mathcal S}|S|(n-|S|)t_S,\\
\text{subject to}\quad &
\sum_{S\in\mathcal S}t_S\delta_S(i,j)\leq d(i,j)
\quad (i<j),\\
&t_S\geq 0\quad(S\in\mathcal S).
\end{aligned}
\end{equation}
It is feasible by taking every $t_S=0$.  For any feasible $t$, the map
$F_t(i)=(t_S\mathbf 1_S(i))_{S\in\mathcal S}$ takes values in
$\ell_1^{|\mathcal S|}$ and is a contraction for $d$.  Its sum of
distances over unordered pairs is exactly the objective in
\eqref{red:primal}.  Consequently \eqref{red:averages} bounds the primal
optimum by $n^2b/2$.  The primal optimum is attained: for each
$S\in\mathcal S$ and any pair separated by $S$, feasibility implies
$0\leq t_S\leq d(i,j)$, so its feasible region is closed and bounded.

The dual of \eqref{red:primal}, with one nonnegative variable $c_{ij}$
for each unordered pair, is
\begin{equation}
\label{red:dual}
\begin{aligned}
\text{minimize}\quad &\sum_{i<j}c_{ij}d(i,j),\\
\text{subject to}\quad &
\sum_{i\in S,\ j\notin S}c_{ij}\geq |S|(n-|S|)
\quad(S\in\mathcal S),\\
&c_{ij}\geq 0\quad(i<j).
\end{aligned}
\end{equation}
Here the cut sum counts each separated unordered pair once.
The dual is feasible, for example with $c_{ij}=1$ for every pair.
Finite-dimensional linear programming duality therefore gives an
optimal dual solution $C$ satisfying
\begin{equation}
\label{red:dual-objective}
\sum_{i<j}c_{ij}d(i,j)\leq \frac{n^2b}{2}.
\end{equation}
Its cut constraints give $\OPT(C)\geq1$ with exactly unit pair demands.

Set $D=\sum_{i<j}d(i,j)\geq n^2a/2>0$.  If
$d(i,j)=|z_i-z_j|^2$ in a Hilbert space, the vectors
$z_i/\sqrt D$ give the feasible normalized SDP distances $d(i,j)/D$.
The span of the finitely many $z_i$ has dimension at most $n$, so these
vectors can be realized in $\R^n$.  The triangles are preserved by
normalization, and \eqref{red:dual-objective} gives
$\GL(C)\leq b/a$.

Finally, the graph of pairs with $c_{ij}>0$ is connected: otherwise the
vertex set of one component would violate a cut constraint in
\eqref{red:dual}.  Choose a spanning tree $\mathcal T$ in this graph
and put $\gamma=\min_{\{i,j\}\in\mathcal T}c_{ij}>0$.  For any feasible
SDP distance array $d'$, the triangle inequalities along tree paths
imply
\[
1=\sum_{i<j}d'(i,j)
\leq \binom n2\sum_{\{i,j\}\in\mathcal T}d'(i,j).
\]
Hence its objective is at least $\gamma/\binom n2>0$.  This proves
$\GL(C)>0$, including when the original semimetric $d$ has zero
distances between distinct vertices.
\end{proof}

\subsection{Replacing the parameter distribution by multiplicities}

Fix a sufficiently large integer $m$, and use the vertex set $V$,
semimetric $d$, and chart map $v_1$ from \Cref{prop:metric}.  Define the
probability weights
\[
w_v=\Pr_{\theta\ \mathrm{uniform\ on}\ T}
\big[v_1(\theta)=v\big],
\qquad T=(-1,1)^m.
\]
Weights on vertices outside the first chart are thus zero.
By \Cref{prop:metric,prop:contraction}, for absolute positive constants
that may change from line to line,
\begin{equation}
\label{red:weighted-bounds}
\begin{gathered}
|V|\leq \exp(Cm\log m),\qquad
\operatorname{diam}(V,d)\leq Cm,\qquad
\sum_{v,z\in V}w_vw_zd(v,z)\geq cm,\\
\sum_{v,z\in V}w_vw_z\|F(v)-F(z)\|_1
\leq C\sqrt m\,(\log m)^{5/2}
\end{gathered}
\end{equation}
for every finite-dimensional $\ell_1$-valued contraction $F$ for $d$.

Every chart vertex must remain present, even if its weight is zero:
its distance constraints still restrict contractions on the first chart.
We therefore give each vertex at least one copy. In the choice below,
the term $m^2|V|$ controls the measure error, while $m^m$ ensures the
lower bound on $\log n$ needed in the final parameter conversion. Set
\[
\begin{gathered}
M=m^m+m^2|V|,
\qquad k_v=1+\lfloor Mw_v\rfloor,\\
\widetilde V=\{(v,h):v\in V,\ 1\leq h\leq k_v\},
\qquad n=|\widetilde V|.
\end{gathered}
\]
Give these copies the semimetric
$\widetilde d((v,h),(z,k))=d(v,z)$.
It is again a squared Hilbert distance satisfying all triangles, by
assigning the same Hilbert vector to every copy of a vertex.  In
particular, distances between copies of one vertex are zero.

Let $\mu_v=k_v/n$ be the pushforward to $V$ of the uniform probability
measure on $\widetilde V$.  To bound the rounding error, write
$e_v=k_v-Mw_v\in(0,1]$ and $E_0=\sum_v e_v$.  Then
$n=M+E_0$, where $0<E_0\leq |V|$, and
\[
\mu_v-w_v=\frac{e_v-E_0w_v}{n}.
\]
Using the convention
$\|\mu-w\|_{\mathrm{TV}}=\frac12\sum_v|\mu_v-w_v|$, we obtain
\begin{equation}
\label{red:tv}
\|\mu-w\|_{\mathrm{TV}}
\leq \frac{E_0}{n}
\leq \frac{|V|}{M}
\leq \frac{1}{m^2},
\qquad
\|\mu\otimes\mu-w\otimes w\|_{\mathrm{TV}}
\leq \frac{2}{m^2}.
\end{equation}
The second inequality follows by writing the difference of the product
measures as
$(\mu-w)\otimes\mu+w\otimes(\mu-w)$ and applying the triangle
inequality for total variation.

For any function with values in $[0,Cm]$, the change in its expectation
under these two product measures is at most $2C/m$.
Applying this to $d$ gives
\begin{equation}
\label{red:uniform-metric}
\frac{1}{n^2}\sum_{x,y\in\widetilde V}\widetilde d(x,y)
=\sum_{v,z\in V}\mu_v\mu_zd(v,z)
\geq c'm
\end{equation}
for all sufficiently large $m$.
Every contraction $\widetilde F$ from $(\widetilde V,\widetilde d)$
into a finite-dimensional $\ell_1$ space agrees on copies of a vertex,
because the distance between those copies is zero.  Since $k_v\geq1$
for every $v$, it therefore induces a contraction $F$ on all of $V$.
The function $\|F(v)-F(z)\|_1$ is bounded by $d(v,z)\leq Cm$.
Thus \eqref{red:weighted-bounds} and \eqref{red:tv} give
\begin{equation}
\label{red:uniform-contraction}
\frac{1}{n^2}\sum_{x,y\in\widetilde V}
\|\widetilde F(x)-\widetilde F(y)\|_1
\leq C'\sqrt m\,(\log m)^{5/2}.
\end{equation}

Apply \Cref{lem:duality} to
\eqref{red:uniform-metric}--\eqref{red:uniform-contraction}.
After labeling $\widetilde V$ by $[n]$, this gives nonnegative
capacities $C^{(m)}$ with unit demands on all distinct pairs and
\begin{equation}
\label{red:m-gap}
\frac{\OPT(C^{(m)})}{\GL(C^{(m)})}
\geq c''\frac{\sqrt m}{(\log m)^{5/2}},
\qquad \GL(C^{(m)})>0.
\end{equation}

\subsection{The number of vertices}

Since $M\leq n\leq M+|V|$, \eqref{red:weighted-bounds} implies, for
an absolute constant $K$ and all sufficiently large $m$,
\begin{equation}
\label{red:size}
m\log m\leq\log n\leq K m\log m.
\end{equation}
In particular $n\to\infty$ and $\log\log n\geq\log m$ for large $m$.
Consequently
\[
\frac{\sqrt{\log n}}{(\log\log n)^3}
\leq \sqrt K\,
\frac{\sqrt m}{(\log m)^{5/2}}.
\]
Combining this with \eqref{red:m-gap} proves \Cref{thm:main}.


\appendix
\section{Gaussian directions: the probabilistic proof}
\label{app:directions}

\begin{proof}[Proof of \Cref{lem:directions}]
Choose all the $g_{s,i}$ independently with the standard Gaussian distribution
on $\R^m$.  We show that the required events hold simultaneously with
probability tending to one as $m\to\infty$.

We first record the elementary estimates used below.  For a standard Gaussian
$G$ in $\R^m$,
\begin{equation}\label{dir:gaussian-identities}
 \E\cos\langle G,w\rangle=e^{-|w|^2/2},
 \qquad
 \E\bigl[G\sin\langle G,w\rangle\bigr]=w e^{-|w|^2/2}.
\end{equation}
Indeed, rotation reduces the first identity to the one-dimensional cosine
transform.  Gaussian integration by parts gives the differential equation
$\phi'(t)=-t\phi(t)$ for that transform, with $\phi(0)=1$.  Differentiating
the first identity with respect to $w$ gives the second; the finite first
Gaussian moment justifies this differentiation.  Also, exponential Markov
inequalities applied to
$\E e^{t|G|^2}=(1-2t)^{-m/2}$, for $t<1/2$, give
for an absolute $c>0$
\begin{equation}\label{dir:gaussian-norm-tail}
 \Pr\left\{|G|\notin[\sqrt m/2,2\sqrt m]\right\}
 \le 2e^{-c m}.
\end{equation}

We will use the following bounded-variable estimate, a form of
Hoeffding's inequality~\cite{Hoe63}. If $X_1,\ldots,X_n$
are independent real random variables with $|X_j|\le B$, then, for
$0<u\le B$,
\begin{equation}\label{dir:bounded-concentration}
 \Pr\left\{\left|\frac1n\sum_{j=1}^n(X_j-\E X_j)\right|>u\right\}
 \le 2\exp\left(-\frac{c n u^2}{B^2}\right).
\end{equation}
For completeness, Taylor expansion, the vanishing first centered moment,
and $|X_j-\E X_j|\le2B$ give
$\E e^{t(X_j-\E X_j)}\le e^{C t^2B^2}$ when $|t|\le1/(2B)$.
Independence and exponential Markov, with $t$ a sufficiently small constant
multiple of $u/B^2$, give both tails in
\eqref{dir:bounded-concentration}.

To obtain \eqref{dir:gaussian-approximation}, truncate each Gaussian summand
on the event $|G|>2\sqrt m$.  By Cauchy--Schwarz and
\eqref{dir:gaussian-norm-tail}, uniformly in $w$,
\begin{equation}\label{dir:truncation-bias}
 \left|\E\bigl[G\sin\langle G,w\rangle
                   \mathbf 1_{\{|G|\le2\sqrt m\}}\bigr]
               -w e^{-|w|^2/2}\right|
 \le \E\bigl[|G|\mathbf 1_{\{|G|>2\sqrt m\}}\bigr]
 \le e^{-c' m},
\end{equation}
after decreasing the positive constant $c'$ and taking $m$ large.
For a fixed unit vector $v$ and fixed $w$, the scalar truncated summand
\[
 \langle v,G\rangle\sin\langle G,w\rangle
       \mathbf 1_{\{|G|\le2\sqrt m\}}
\]
has absolute value at most $2\sqrt m$.  Thus
\eqref{dir:bounded-concentration}, with $n=N$ and
$u=1/(2\sqrt m)$, bounds its empirical-mean deviation probability by
$2e^{-cN/m^2}$.

Choose a $1/8$-net $\mathcal V$ of the unit sphere and an $m^{-3}$-net
$\mathcal W$ of the ball of radius $m^{42}$.  The usual disjoint-ball packing
argument gives
\[
 |\mathcal V|\le17^m,
 \qquad
 |\mathcal W|\le(1+2m^{45})^m
                 =\exp(O(m\log m)).
\]
Union bound the preceding scalar deviations over $s$, $v\in\mathcal V$,
and $w\in\mathcal W$.  The failure probability tends to zero, since
$N/m^2=m^4$.  On the resulting event, for every $s$ and $w\in\mathcal W$,
the vector deviation from its truncated expectation is at most
$4/(7\sqrt m)$: bounds on inner products with a $1/8$-net control the
Euclidean norm with factor $8/7$.

Each truncated empirical vector function is $4m$-Lipschitz in $w$, because
the operator norm of its derivative is at most the average of the
truncated $|g_{s,i}|^2$.  The function $w\mapsto we^{-|w|^2/2}$ is
$2$-Lipschitz.  Consequently the bias bound
\eqref{dir:truncation-bias} and interpolation from $\mathcal W$ give,
uniformly on the whole ball, an error at most
\[
 \frac{4}{7\sqrt m}+e^{-c'm}+(4m+2)m^{-3}
 \le \frac2{\sqrt m}
\]
for all sufficiently large $m$.  This proves the desired approximation
for the truncated samples.

The covariance estimate follows from the same argument with quadratic
forms.  For a fixed unit vector $v$, the variable
$\langle v,G\rangle^2\mathbf 1_{\{|G|\le2\sqrt m\}}$ is bounded by $4m$,
and its expectation differs from $1$ by at most
\[
 \bigl(\E\langle v,G\rangle^4\bigr)^{1/2}
 \Pr\{|G|>2\sqrt m\}^{1/2}
 \le e^{-c'm},
\]
where constants can again be adjusted, since
$\E\langle v,G\rangle^4=3$.  Apply
\eqref{dir:bounded-concentration} with $u=1/8$, and union bound over $s$
and $v\in\mathcal V$.  With probability tending to one, every truncated
empirical covariance matrix $M_s$ satisfies
\[
 \|M_s-I_m\|_{\op}
 \le\frac{1}{1-2/8}\left(\frac18+e^{-c'm}\right)
 <\frac12.
\]
Here we used the elementary symmetric-matrix net estimate
$\|A\|_{\op}\le(1-2\rho)^{-1}\max_{v\in\mathcal V}|\langle Av,v\rangle|$
for a sphere net of radius $\rho<1/2$.

By \eqref{dir:gaussian-norm-tail}, all $SN=m^9$ original Gaussian samples
satisfy \eqref{dir:norms} with probability tending to one.  On this event
the truncations change no sample, so the preceding conclusions establish
the first two assertions.  Pairwise nonparallelism within each family holds
almost surely.

It remains to establish \eqref{dir:good-charts}.  For a fixed unit vector
$v$, the one-dimensional Gaussian tail bound gives
\[
 \Pr\left\{\max_i|\langle v,g_{s,i}\rangle|
                      >10\sqrt{\log m}\right\}
 \le 2N e^{-50\log m}=2m^{-44}.
\]
Call such an index $s$ bad for $v$.  Before any conditioning, these bad-index
events are independent across $s$.  Therefore the probability that at least
$S/2$ indices are bad for this fixed $v$ is at most
\[
 2^{S}(2m^{-44})^{S/2}.
\]
Take a sphere net of radius $1/m$, with cardinality at most $(1+2m)^m$,
and union bound this last estimate over its points.  The resulting failure
probability tends to zero.  Finally intersect with the already established
norm event.  If $v_0$ is a net point with $|v-v_0|\le1/m$, then
\[
 |\langle v-v_0,g_{s,i}\rangle|\le\frac2{\sqrt m}
\]
for every $s,i$.  Every chart good for $v_0$ consequently satisfies
\eqref{dir:good-charts} for $v$, with $C=11$ for all sufficiently large
$m$.  This argument uses independence only for the unconditioned Gaussian
samples.  A final union bound intersects this event with the preceding
high-probability events and proves the lemma.
\end{proof}


\bibliographystyle{plain}
\bibliography{references}
\end{document}
